% !PRINTED 466
\weberpage{466}

%<<CONT
denote by $\Omega(\varrho)$. If $\mu$ is the degree of $N$, then $\varrho$ satisfies an irreducible equation of degree $\mu$
\weq{5}{g(t) = 0,}
of which it is to be shown that it is a normal equation. For this purpose we first remark that the one root $\varrho$ of this equation is a rational function of $\alpha, \alpha_1 \ldots \alpha_{m-1}$, because it was contained in $N$. To indicate this, let us set
\[
\varrho = \varrho(\alpha, \alpha_1 \ldots \alpha_{m-1}),
\]
and now form all the different arrangements of the digits
\[
0,\ 1,\ 2 \ldots m - 1,
\]
whose number amounts to $\Pi(m)$:
\weq{6}{(0, 1, 2 \ldots m - 1),\ (0', 1', 2' \ldots (m-1)'),\ (0'', 1'', 2'' \ldots (m-1)'')\,\ldots,}
in which the digits with one, two, etc. accents are the same as those without accent, only in a different order, and from these form the functions
\weq{7}{\varrho = \varrho(\alpha_0, \alpha_1 \ldots \alpha_{m-1}),\ \varrho' = \varrho(\alpha_{0'}, \alpha_{1'} \ldots \alpha_{(m-1)'}),\ \varrho'' = \varrho(\alpha_{0''}, \alpha_{1''} \ldots \alpha_{(m-1)''})\,\ldots,}
without concern as to whether some among them may be equal to one another or not.

If in all the arrangements (6) we carry out one and the same permutation, e.g. 0 with 1, the totality of these arrangements does not change; only their order becomes a different one. For, first, such a permutation can produce nothing other than arrangements of the digits, and secondly, two different arrangements cannot pass into the same arrangement by one and the same permutation.

Hence by every such permutation the functions (7) will likewise only be permuted among one another. If, then, we form the product
\[
G(t) = (t - \varrho)\,(t - \varrho')\,(t - \varrho'')\,\ldots
\]
for a variable $t$, its coefficients, which are certainly functions of $\alpha, \alpha_1 \ldots \alpha_{m-1}$, remain unchanged when these quantities are permuted in any way; that is, they are symmetric functions of the roots of the function $F(x)$, and hence $G(t)$ is a function of $t$ in $\Omega$. All the roots of $G(t)$ are quantities in $N$, since they are rationally expressed by the $\alpha$.
